% Sections 2.1 to 2.7, translated from sv/chapters/02/theory.tex.

\section{Variables and algebraic expressions}
\label{c2:sec:variabler}
A \textbf{variable} is an unknown number represented by a letter, for example $x$, $R$, $U$ or
$I$.

An \textbf{algebraic expression} is a combination of variables, constants and operations, for
example
\[
  3x + 2y - 5, \qquad \frac{U}{R}, \qquad 2R_1 + R_2.
\]

\subsection{Terms, factors and coefficients}
\label{c2:sec:termer}
In the expression $5x^2 - 3x + 7$:
\begin{itemize}
  \item $5x^2$, $-3x$ and $7$ are \textbf{terms}. The terms are separated by $+$ or $-$.
  \item $5$ is the \textbf{coefficient} of $x^2$ and $-3$ is the coefficient of $x$, while $7$ is
    a \textbf{constant}.
  \item In the product $5 \cdot x^2$, $5$ and $x^2$ are \textbf{factors}.
\end{itemize}

\section{Simplifying: collecting like terms}
\label{c2:sec:forenkling}
Terms with the \emph{same variable part} are called \textbf{like terms} and can be added:
\[
  3x + 5x = 8x, \qquad 4R - R = 3R, \qquad 2x^2 + 3x - x^2 + x = x^2 + 4x
\]
Terms with \emph{different} variable parts can\textbf{not} be simplified further:
\[
  3x + 2y \quad \text{(cannot be simplified further)}
\]

\section{The distributive law}
\label{c2:sec:distributiv}
The distributive law is used to expand brackets:
\[
  a(b + c) = ab + ac
\]

\begin{exempel}
\begin{gather*}
  3(2x + 4) = 6x + 12 \\
  -2(R_1 + R_2 - 3) = -2R_1 - 2R_2 + 6
\end{gather*}
\end{exempel}

The distributive law is used in reverse when factorising; see \secref{c2:sec:faktorsattning}.

\section{The squaring rules}
\label{c2:sec:kvadrering}
The squaring rules are used a great deal in electrical engineering, for example in calculating
power and impedance.

\subsection{Square of a sum}
\label{c2:sec:kvsumma}
\[
  (a + b)^2 = a^2 + 2ab + b^2
\]

\begin{exempel}
\[
  (R + 2)^2 = R^2 + 4R + 4
\]
\end{exempel}

\subsection{Square of a difference}
\label{c2:sec:kvdiff}
\[
  (a - b)^2 = a^2 - 2ab + b^2
\]

\begin{exempel}
\[
  (U - 3)^2 = U^2 - 6U + 9
\]
\end{exempel}

\section{The difference of two squares}
\label{c2:sec:konjugat}
\[
  (a + b)(a - b) = a^2 - b^2
\]

\begin{exempel}
\begin{gather*}
  (x + 5)(x - 5) = x^2 - 25 \\
  (2R + 1)(2R - 1) = 4R^2 - 1
\end{gather*}
\end{exempel}

\section{Factorisation}
\label{c2:sec:faktorsattning}
Factorisation means writing an expression as a \textbf{product of factors}, the reverse of
expanding. There are several methods.

\subsection{Common factor}
\label{c2:sec:gemensam}
Identify the factor that all the terms have in common and take it outside a bracket:
\begin{gather*}
  6x + 9 = 3(2x + 3) \\
  4R_1^2 + 8R_1 = 4R_1(R_1 + 2) \\
  RI^2 - RI = RI(I - 1)
\end{gather*}

\subsection{The difference of two squares in reverse}
\label{c2:sec:konjbakat}
An expression of the form $a^2 - b^2$ can be factorised:
\[
  a^2 - b^2 = (a + b)(a - b)
\]

\begin{exempel}
\begin{gather*}
  x^2 - 16 = (x + 4)(x - 4) \\
  4U^2 - 9 = (2U + 3)(2U - 3)
\end{gather*}
\end{exempel}

\subsection{The squaring rules in reverse}
\label{c2:sec:kvbakat}
Expressions of the form $a^2 + 2ab + b^2$ or $a^2 - 2ab + b^2$ can be factorised:
\begin{gather*}
  x^2 + 6x + 9 = (x + 3)^2 \\
  R^2 - 4R + 4 = (R - 2)^2
\end{gather*}

\section{Application in electrical engineering}
\label{c2:sec:tillampning}
\textbf{Ohm's law} in algebraic form is
\[
  U = R \cdot I.
\]
The law can be rewritten depending on which quantity is sought:
\[
  R = \frac{U}{I}, \qquad I = \frac{U}{R}
\]
The \textbf{power formula} gives rise to algebraic expressions:
\[
  P = U \cdot I = R \cdot I^2 = \frac{U^2}{R}
\]
For the resistors $R_1$ and $R_2$ in series,
\[
  R_{\text{TOT}} = R_1 + R_2,
\]
which is a linear algebraic expression.
